\section{薄板的谱与振动}
\label{chap:phys-sm-thin-plate-chladni}

薄板方程已经把三维弹性约化为二维四阶边值问题。真正可见的图样却由边界条件和本征函数决定：同一块板怎样选出共振频率，节线又为什么正好落在本征函数的零集上？

\subsection{圆形板}
\label{sec:phys-sm-eigenvalue}

前一节得到的变分方程允许弯曲刚度 \(\mathsf D(\vb r)\) 和面质量 \(\mu_A(\vb r)\) 随位置改变。为了把本征值实际算出来，现在取均匀各向同性板，使弯曲刚度成为常数 \(D\)、面质量成为常数 \(\rho_{\rm m} h\)，并把中面区域取为环域
\[
 \mathcal A=\{(r,\theta):R_0<r<R\},
 \qquad
 \xi=\frac rR,
 \qquad
 \xi_0=\frac{R_0}{R}.
\]
这一节的目的不是为每一种边界分别建立一套公式，而是把同一个自伴四阶算子在圆对称坐标中的谱、Green 算子和受迫响应写成直接联系。

自由振动取 \(f_\perp=0\)，并令
\[
 w(r,\theta,t)=\Re\{\Psi(r,\theta)e^{-\ii\omega t}\}.
\]
均匀各向同性板的空间本征方程为
\begin{equation}
 (\nabla_t^4-k^4)\Psi=0,
 \qquad
 k^4=\frac{\rho_{\rm m} h\omega^2}{D}.
 \label{eq:phys-sm-annulus-eigenproblem}
\end{equation}
圆对称使角向 Fourier 扇区彼此独立。写
\[
 \Psi(r,\theta)=\mathcal R_m(r)e^{\ii m\theta},
 \qquad m\in\mathbb Z,
\]
并定义 \(\beta=kR\)。下式中的 \(J_m,Y_m\) 是第一、第二类 Bessel 函数，\(I_m,K_m\) 是相应的修正 Bessel 函数。径向齐次解空间由
\begin{equation}
 \vb b_m(\xi;\beta)
 :=\bigl(J_m(\beta\xi),\,Y_m(\beta\xi),\,I_m(\beta\xi),\,K_m(\beta\xi)\bigr)
 \label{eq:phys-sm-annulus-radial-basis}
\end{equation}
张成。于是 \(\mathcal R_m=\vb b_m\vb c\)，其中 \(\vb c\in\mathbb C^4\) 是待定系数列向量。

边界条件最好直接用变分共轭迹表示，而不是把 Bessel 恒等式提前展开。对 \(\mathcal R_m(r)e^{\ii m\theta}\)，定义
\begin{align*}
 M_{rr}[\mathcal R_m]
 &=-D\left[\mathcal R_m''
 +\nu\left(\frac1r\mathcal R_m'-\frac{m^2}{r^2}\mathcal R_m\right)\right],\\
 V_r[\mathcal R_m]
 &=-D\left[\mathcal R_m'''
 +\frac1r\mathcal R_m''
 -\frac{1+(2-\nu)m^2}{r^2}\mathcal R_m'
 +\frac{(3-\nu)m^2}{r^3}\mathcal R_m\right].
\end{align*}
三类经典边界迹因此统一为
\[
 \mathcal B_{\mathrm C}[\mathcal R]=(\mathcal R,\mathcal R'),
 \qquad
 \mathcal B_{\mathrm S}[\mathcal R]=(\mathcal R,M_{rr}),
 \qquad
 \mathcal B_{\mathrm F}[\mathcal R]=(M_{rr},V_r).
\]
若外缘类型为 \(X\in\{\mathrm C,\mathrm S,\mathrm F\}\)，内缘类型为 \(Y\in\{\mathrm C,\mathrm S,\mathrm F\}\)，则九种有序边界组合全部由同一个 \(4\times4\) 矩阵表示：
\begin{equation}
 \vb Z_m^{XY}(\beta,\xi_0)
 :=\begin{pmatrix}
 \mathcal B_X[\vb b_m](R)\\
 \mathcal B_Y[\vb b_m](R_0)
 \end{pmatrix},
 \qquad
 \det\vb Z_m^{XY}(\beta,\xi_0)=0.
 \label{eq:phys-sm-annulus-characteristic}
\end{equation}
这里每个 \(\mathcal B\) 产生两行，作用在行向量 \(\vb b_m\) 的每个分量上；当边界算子作用于 \(\vb b_m(\xi;\beta)\) 时，径向导数均按 \(\partial_r=R^{-1}\partial_\xi\) 还原为物理导数。内缘外法向与 \(+\vb e_r\) 相反，但齐次条件中的整体符号不影响零空间。\eqnrefs{eq:phys-sm-annulus-characteristic} 的正根按从小到大记为 \(\beta_{m,s}^{XY}\)，其中 \(s=1,2,\ldots\) 是径向根序号，于是
\begin{equation}
 \omega_{m,s}^{XY}
 =\frac{(\beta_{m,s}^{XY})^2}{R^2}
 \sqrt{\frac{D}{\rho_{\rm m} h}}.
 \label{eq:phys-sm-annulus-frequency}
\end{equation}
实心圆板只是 \(\xi_0\to0\) 并要求中心正则的极限，此时 \(Y_m\) 与 \(K_m\) 分支自动排除。

对 \(m\ge1\)，实本征函数可取 \(\mathcal R_{m,s}\cos m\theta\) 与 \(\mathcal R_{m,s}\sin m\theta\)，二者具有相同频率。任意实线性组合只会旋转这组节径；理想圆对称下其方向并不唯一，实验中的微弱几何或驱动不对称才会选定具体组合。径向零点给出节圆，因此 Chladni 图样并不是额外模型，而是被实际激发本征子空间中某个线性组合的节点集。用单一指标 \(j\) 统摄 \((m,s)\) 与角向简并指标，并定义质量加权模态范数
\[
 \mathcal N_j
 :=\int_{\mathcal A}\rho_{\rm m} h\,|\Psi_j|^2\dd a.
\]
这里的 \(j\) 只是离散模态标签，不与材料的 Poisson 比 \(\nu\) 混用。一般载荷下的模态坐标满足
\begin{equation}
 \ddot W_j+\omega_j^2W_j
 =\frac{1}{\mathcal N_j}
 \int_{\mathcal A}f_\perp(\vb r,t)\Psi_j^*(\vb r)\dd a.
 \label{eq:phys-sm-annulus-modal-dynamics}
\end{equation}

\label{sec:phys-sm-green-dynamic}

谱分解立即把静态、点载和谐响应放进同一个预解式。若所选边界实现没有零模，静态 Green 核满足
\[
 D\nabla_t^4G^{XY}(\vb r,\vb r_*)
 =\delta^{(2)}(\vb r-\vb r_*),
\]
并对第一个变量满足 \(XY\) 边界条件。由\eqnrefs{eq:phys-sm-annulus-modal-dynamics} 的本征系可写成
\begin{equation}
 G^{XY}(\vb r,\vb r_*)
 =\sum_j
 \frac{\Psi_j(\vb r)\Psi_j^*(\vb r_*)}
 {\mathcal N_j\omega_j^2}.
 \label{eq:phys-sm-annulus-static-green-spectral}
\end{equation}
因此任意静载都由
\[
 w(\vb r)=\int_{\mathcal A}G^{XY}(\vb r,\vb r_*)f_\perp(\vb r_*)\dd a_*
\]
给出，集中力只是在积分中取 \(\delta\) 源。对驱动角频率为 \(\omega_{\rm d}\) 的谐载，避开无阻尼共振极点时相应动态核为
\begin{equation}
 G_{\omega_{\rm d}}^{XY}(\vb r,\vb r_*)
 =\sum_j
 \frac{\Psi_j(\vb r)\Psi_j^*(\vb r_*)}
 {\mathcal N_j(\omega_j^2-\omega_{\rm d}^2)}.
 \label{eq:phys-sm-annulus-dynamic-green-spectral}
\end{equation}
这两个式子分别是前一章 \(\mathsf L_X^{-1}\) 与动态预解式的环板实现；几何只决定本征根与本征函数，不改变 Green 理论本身。

若两条边都是自由边界，完整二维问题含 \(1,x,y\) 三个刚体零模，静态普通逆不存在。取关于质量内积
\[
 \langle f,g\rangle_\mu
 :=\int_{\mathcal A}\rho_{\rm m} h\,f^*g\,\dd a
\]
相互正交的刚体基 \(\{\Psi_{0,\ell}\}_{\ell=1}^3\)，并记
\(
 \mathcal N_{0,\ell}=\langle\Psi_{0,\ell},\Psi_{0,\ell}\rangle_\mu
\)。严格静态问题必须把这三个模投影掉并使用 Moore--Penrose 逆；这正对应净横向力和两个倾覆力矩的兼容条件。可是对 \(\omega_{\rm d}\neq0\) 的谐驱动，刚体模并没有消失，而是给出真实的整体惯性响应：
\begin{equation}
 G_{\omega_{\rm d}}^{\mathrm{FF}}(\vb r,\vb r_*)
 =-\sum_{\ell=1}^{3}
 \frac{\Psi_{0,\ell}(\vb r)\Psi_{0,\ell}^*(\vb r_*)}
 {\mathcal N_{0,\ell}\omega_{\rm d}^2}
 +\sum_{j:\,\omega_j>0}
 \frac{\Psi_j(\vb r)\Psi_j^*(\vb r_*)}
 {\mathcal N_j(\omega_j^2-\omega_{\rm d}^2)}.
 \label{eq:phys-sm-annulus-ff-dynamic-zero-modes}
\end{equation}
因此“投影掉零模”只属于静态伪逆，不属于一般动态预解式。轴对称扇区只保留常数刚体模，下面可把静态投影写得完全显式。

\label{sec:phys-sm-eight-green}

取 \(m=0\)。这里的轴对称 Green 核按环向平均测度 \(2\pi\xi\,\dd\xi\) 归一化，因此源项 \(\delta(\xi-\eta)/(2\pi\eta)\) 表示总力为一的均匀环载，而不是固定方位角处的二维点力。若物理环载总力为 \(P\)、作用半径为 \(r_*=R\eta\)，则
\[
 f_\perp(r)=\frac{P}{2\pi r_*}\delta(r-r_*),
 \qquad
 w(r)=\frac{PR^2}{D}\,\mathcal G(\xi,\eta).
\]
无量纲轴对称算子和两个自然边界迹定义为
\[
 \mathcal L_0
 :=\left(\frac{\dd^2}{\dd\xi^2}
 +\frac1\xi\frac{\dd}{\dd\xi}\right)^2,
 \qquad
 \mathcal M[f]:=f''+\frac{\nu}{\xi}f',
 \qquad
 \mathcal V[f]:=f'''+\frac1\xi f''-\frac1{\xi^2}f'.
\]
物理量分别满足 \(M_{rr}=-(D/R^2)\mathcal M\)、\(V_r=-(D/R^3)\mathcal V\)。对无零模的边界组合，令 \(\{u_1^Y,u_2^Y\}\) 为自动满足内缘 \(Y\) 条件的齐次基，\(\{v_1^X,v_2^X\}\) 为自动满足外缘 \(X\) 条件的齐次基，则
\begin{equation}
 \mathcal G_{XY}(\xi,\eta)=
 \begin{cases}
  c_{-,1}\,u_1^Y(\xi)+c_{-,2}\,u_2^Y(\xi), & \xi_0<\xi<\eta,\\[2pt]
  c_{+,1}\,v_1^X(\xi)+c_{+,2}\,v_2^X(\xi), & \eta<\xi<1,
 \end{cases}
 \label{eq:phys-sm-axisym-green-piecewise}
\end{equation}
并满足
\[
 \mathcal L_0\mathcal G_{XY}(\xi,\eta)
 =\frac{\delta(\xi-\eta)}{2\pi\eta}.
\]
位移、斜率和弯矩连续，只有有效剪力跳跃：
\begin{equation}
 \mathcal V[\mathcal G](\eta^+)
 -\mathcal V[\mathcal G](\eta^-)
 =\frac{1}{2\pi\eta}.
 \label{eq:phys-sm-axisym-green-jump}
\end{equation}
于是令
\[
 \vb c_{XY}:=(c_{-,1},c_{-,2},c_{+,1},c_{+,2})^{\mathsf T},
\]
四个局域系数由统一界面系统
\begin{equation}
 \vb H_{XY}(\eta)\vb c_{XY}(\eta)
 =\begin{pmatrix}0\\0\\0\\(2\pi\eta)^{-1}\end{pmatrix},
 \label{eq:phys-sm-green-interface-system}
\end{equation}
决定，其中
\begin{equation}
\vb H_{XY}(\eta)=
\left.
\begin{pmatrix}
 u_1^Y & u_2^Y & -v_1^X & -v_2^X\\
 (u_1^Y)' & (u_2^Y)' & -(v_1^X)' & -(v_2^X)'\\
 \mathcal M[u_1^Y] & \mathcal M[u_2^Y] & -\mathcal M[v_1^X] & -\mathcal M[v_2^X]\\
 -\mathcal V[u_1^Y] & -\mathcal V[u_2^Y] & \mathcal V[v_1^X] & \mathcal V[v_2^X]
\end{pmatrix}
\right|_{\xi=\eta}.
\label{eq:phys-sm-green-interface-matrix}
\end{equation}
因此所有轴对称 \(\mathrm C/\mathrm S/\mathrm F\) 组合只是在改变两端二维齐次子空间；源点跳跃是同一个四阶 Green 结构。
自由--自由情形中，轴对称扇区只剩常数零模。Moore--Penrose 核必须把 \(\delta\) 源投影到常数的正交补，而不是任意指定某一点的位移：
\begin{equation}
 \mathcal L_0\mathcal G_{\mathrm{FF}}^{+}(\xi,\eta)
 =\frac{\delta(\xi-\eta)}{2\pi\eta}
 -\frac{1}{\pi(1-\xi_0^2)},
 \qquad
 \int_{\xi_0}^{1}2\pi\xi\,
 \mathcal G_{\mathrm{FF}}^{+}(\xi,\eta)\dd\xi=0.
 \label{eq:phys-sm-axisym-ff-pseudoinverse}
\end{equation}
这正是\chapref{chap:phys-cont-biharmonic}自由边界伪逆的轴对称实现。
\begin{exbox}[外缘固支、内缘自由环板的均匀横载]
一般 Green 理论建立以后，均匀载荷才作为特例出现。令 \(f_\perp(r)=f_0\)，外缘 \(\xi=1\) 固支、内缘 \(\xi=\xi_0\) 自由。轴对称静力方程为 \(\mathcal L_0f=1\)，且
\[
 w(r)=\frac{f_0R^4}{D}f(\xi).
\]
定义
\[
 \Delta=(1-\nu)+(1+\nu)\xi_0^2.
\]
满足 \(f(1)=f'(1)=0\)、\(\mathcal M[f](\xi_0)=\mathcal V[f](\xi_0)=0\) 的解为
\begin{equation}
 f(\xi)=\frac{\xi^4}{64}
 -\frac{\xi_0^2}{8}\xi^2\ln\xi
 +b_{\mathrm{CF}}\xi^2+c_{\mathrm{CF}}\ln\xi+d_{\mathrm{CF}},
 \label{eq:phys-sm-cf-uniform-deflection}
\end{equation}
其中
\begin{align*}
 b_{\mathrm{CF}}
 &=\frac{(\nu+3)\xi_0^4+2(1-\nu)\xi_0^2+(\nu-1)
 +4(1+\nu)\xi_0^4\ln\xi_0}{32\Delta},\\
 c_{\mathrm{CF}}
 &=\frac{\xi_0^2\bigl[(\nu-1)\xi_0^2-(1+\nu)
 -4(1+\nu)\xi_0^2\ln\xi_0\bigr]}{16\Delta},\\
 d_{\mathrm{CF}}
 &=\frac{(1-\nu)+(3\nu-5)\xi_0^2-2(\nu+3)\xi_0^4
 -8(1+\nu)\xi_0^4\ln\xi_0}{64\Delta}.
\end{align*}
当 \(\xi_0\to0\) 时，对数奇异分支自动消失，得到实心固支圆板
\[
 w(r)=\frac{f_0R^4}{64D}
 \left(1-\frac{r^2}{R^2}\right)^2.
\]
这个极限同时检查了载荷归一化、边界条件和 \(R^4/D\) 的量纲。
\end{exbox}
\begin{proof}[解]
定义径向 Bessel 算子
\begin{equation*}
 \Delta_m:=\frac{\dd^2}{\dd r^2}
 +\frac1r\frac{\dd}{\dd r}-\frac{m^2}{r^2}.
\end{equation*}
由角向本征关系 \(\partial_\theta^2 e^{\ii m\theta}=-m^2 e^{\ii m\theta}\)，面内拉普拉斯算子对分离变量 \(\Psi=\mathcal R_m(r)e^{\ii m\theta}\) 的作用为
\begin{equation*}
 \nabla_t^2\Psi
 =e^{\ii m\theta}
 \left(\mathcal R_m''+\frac1r\mathcal R_m'-\frac{m^2}{r^2}\mathcal R_m\right)
 =e^{\ii m\theta}\Delta_m\mathcal R_m,
\end{equation*}
由于 \(e^{\ii m\theta}\) 仍是 \(\partial_\theta^2\) 的本征函数，再作用一次得 \(\nabla_t^4\Psi=e^{\ii m\theta}\Delta_m^2\mathcal R_m\)。把分离变量代入\eqnrefs{eq:phys-sm-annulus-eigenproblem}，约去公因子 \(e^{\ii m\theta}\) 后得到
\begin{equation*}
 (\Delta_m^2-k^4)\mathcal R_m=0.
\end{equation*}
由于 \(\Delta_m\) 与常数 \(k^2\) 可换，算子多项式可因式分解为
\begin{equation*}
 \Delta_m^2-k^4
 =(\Delta_m-k^2)(\Delta_m+k^2)
 =(\Delta_m+k^2)(\Delta_m-k^2).
\end{equation*}
于是齐次解空间是两个二阶算子核的直和。\(\Delta_m+k^2\) 的核满足
\begin{equation*}
 \mathcal R_m''+\frac1r\mathcal R_m'
 +\left(k^2-\frac{m^2}{r^2}\right)\mathcal R_m=0,
\end{equation*}
两端乘 \(r^2\) 并作变量替换 \(z=kr\)，即化为标准 Bessel 方程 \(z^2 R''+zR'+(z^2-m^2)R=0\)，其两个线性无关解为 \(J_m(kr),Y_m(kr)\)。同理 \(\Delta_m-k^2\) 的核满足把 \(k^2\) 换成 \(-k^2\) 的方程，作同一替换得到修正 Bessel 方程 \(z^2 R''+zR'-(z^2+m^2)R=0\)，两个线性无关解为 \(I_m(kr),K_m(kr)\)。四个解在 \(r>0\) 上线性无关，因此径向齐次解空间即\eqnrefs{eq:phys-sm-annulus-radial-basis}。
设 \(\mathcal R_m=\vb b_m\vb c\)，其中 \(\vb c=(c_1,c_2,c_3,c_4)^{\mathsf T}\)。边界迹 \(\mathcal R\)、\(\mathcal R'\)、\(M_{rr}\)、\(V_r\) 中每一个都是 \(\mathcal R_m\) 的常系数线性微分算子，因而对 \(\vb c\) 线性。具体地，外缘 \(r=R\) 上的边界算子 \(\mathcal B_X\) 作用在基向量 \(\vb b_m\) 上得到一个 \(2\times4\) 行块
\begin{equation*}
 \mathcal B_X[\vb b_m](R)
 =\begin{pmatrix}
  \mathcal B_{X,1}[b_{m,1}](R)&\cdots&\mathcal B_{X,1}[b_{m,4}](R)\\
  \mathcal B_{X,2}[b_{m,1}](R)&\cdots&\mathcal B_{X,2}[b_{m,4}](R)
 \end{pmatrix},
\end{equation*}
内缘 \(r=R_0\) 上同理得到 \(\mathcal B_Y[\vb b_m](R_0)\)。把两个行块上下堆叠，便得到 \(4\times4\) 矩阵 \(\vb Z_m^{XY}\) 和齐次线性系统
\begin{equation*}
 \vb Z_m^{XY}(\beta,\xi_0)\vb c=0.
\end{equation*}
非零 \(\vb c\) 存在当且仅当 \(\vb Z_m^{XY}\) 的秩小于四，即其行列式为零
\begin{equation*}
 \det\vb Z_m^{XY}(\beta,\xi_0)=0.
\end{equation*}
此即\eqnrefs{eq:phys-sm-annulus-characteristic}。因此 \(\mathrm C,\mathrm S,\mathrm F\) 在外缘和内缘的九种有序组合只是 \(X,Y\) 的九个取值，而不是九条独立推导链。
\end{proof}
\begin{proof}[解]
轴对称双调和齐次方程 \(\mathcal L_0 f=0\) 中，\(\mathcal L_0=(\dd_\xi^2+\xi^{-1}\dd_\xi)^2\)。先解一阶因子 \((\dd_\xi^2+\xi^{-1}\dd_\xi)g=0\)，即 \((\xi g')'=0\)，得 \(g=A\ln\xi+B\)；再解 \((\dd_\xi^2+\xi^{-1}\dd_\xi)g=\ln\xi\) 与 \(=1\) 两类非齐次方程，得到基本解空间
\[
 \operatorname{span}\{1,\xi^2,\ln\xi,\xi^2\ln\xi\}.
\]
在该四维空间中施加任一端点的两个边界条件，就得到二维齐次子空间。例如可取内缘基
\[
\begin{aligned}
 u_1^{\mathrm C}&=1-\frac{\xi^2}{\xi_0^2}+2\ln\frac{\xi}{\xi_0},&
 u_2^{\mathrm C}&=1-\frac{\xi^2}{\xi_0^2}
 +2\frac{\xi^2}{\xi_0^2}\ln\frac{\xi}{\xi_0},\\
 u_1^{\mathrm S}&=2(1+\nu)\ln\frac{\xi}{\xi_0}
 +(1-\nu)\left(\frac{\xi^2}{\xi_0^2}-1\right),&
 u_2^{\mathrm S}&=(\nu+3)\left(1-\frac{\xi^2}{\xi_0^2}\right)
 +2(1+\nu)\frac{\xi^2}{\xi_0^2}\ln\frac{\xi}{\xi_0},\\
 u_1^{\mathrm F}&=1,&u_2^{\mathrm F}&=u_1^{\mathrm S}.
\end{aligned}
\]
对应外缘基由 \(\xi_0\to1\) 得到。以 \(\mathrm C\) 内缘为例验证：在 \(\xi=\xi_0\) 处 \(u_1^{\mathrm C}\) 的值为 \(1-1+2\ln1=0\)，导数为 \(-2\xi/\xi_0^2+2/\xi|_{\xi_0}=-2/\xi_0+2/\xi_0=0\)；\(u_2^{\mathrm C}\) 同样有 \(1-1+2\cdot1\cdot\ln1=0\) 与 \(-2/\xi_0+2/\xi_0+4\ln1/\xi_0=0\)。对 \(\mathrm S\) 内缘，利用 \(\mathcal M[f]=f''+\nu(f'+f/\xi)/\xi\)，在 \(\xi=\xi_0\) 处 \(u_1^{\mathrm S}\) 给出 \(f=0\) 且 \(\mathcal M[f]=0\)；\(\mathrm F\) 内缘则由 \(u_1^{\mathrm F}=1\) 给出 \(\mathcal M[1]=0\)、\(\mathcal V[1]=0\)。逐项代入可验证 \(\mathrm C:(f,f')=(0,0)\)、\(\mathrm S:(f,\mathcal M[f])=(0,0)\)、\(\mathrm F:(\mathcal M[f],\mathcal V[f])=(0,0)\)。
另一方面，利用 Kelvin 恒等式可将四阶算子写成散度形式
\[
 \mathcal L_0f
 =\frac1\xi\frac{\dd}{\dd\xi}\!\left[\xi\mathcal V[f]\right],
\]
其中 \(\mathcal V[f]=f'''+\xi^{-1}f''-\xi^{-2}f'\) 是有效剪力算子。把 Green 方程 \(\mathcal L_0\mathcal G=\delta(\xi-\eta)/(2\pi\eta)\) 乘 \(\xi\) 后在 \((\eta-\varepsilon,\eta+\varepsilon)\) 上积分，左端由散度定理给出
\[
 \int_{\eta-\varepsilon}^{\eta+\varepsilon}
 \frac{\dd}{\dd\xi}\!\left[\xi\mathcal V[\mathcal G]\right]\dd\xi
 =\left[\xi\mathcal V[\mathcal G]\right]_{\eta^-}^{\eta^+},
\]
右端为
\[
 \int_{\eta-\varepsilon}^{\eta+\varepsilon}
 \frac{\delta(\xi-\eta)}{2\pi\eta}\,\xi\,\dd\xi
 =\frac1{2\pi}.
\]
令 \(\varepsilon\downarrow0\) 即得到\eqnrefs{eq:phys-sm-axisym-green-jump}。物理上，源点处有效剪力的跳跃对应于作用在环 \(\xi=\eta\) 上的单位环向线载荷；位移、斜率与弯矩连续则保证板在该处没有嵌入裂纹或扭折。
再把分段解\eqnrefs{eq:phys-sm-axisym-green-piecewise} 代入三条连续条件 \(f(\eta^-)=f(\eta^+)\)、\(f'(\eta^-)=f'(\eta^+)\)、\(\mathcal M[f](\eta^-)=\mathcal M[f](\eta^+)\) 与跳跃条件 \(\mathcal V[f](\eta^+)-\mathcal V[f](\eta^-)=1/(2\pi\eta)\)，把四个未知系数 \(\vb c_{XY}=(c_{-,1},c_{-,2},c_{+,1},c_{+,2})^{\mathsf T}\) 排成列向量，逐行得到矩阵\eqnrefs{eq:phys-sm-green-interface-matrix} 和系统\eqnrefs{eq:phys-sm-green-interface-system}。具体地，前两行来自位移与斜率连续：
\[
 c_{-,1}u_1^Y(\eta)+c_{-,2}u_2^Y(\eta)
 =c_{+,1}v_1^X(\eta)+c_{+,2}v_2^X(\eta),
\]
\[
 c_{-,1}(u_1^Y)'(\eta)+c_{-,2}(u_2^Y)'(\eta)
 =c_{+,1}(v_1^X)'(\eta)+c_{+,2}(v_2^X)'(\eta),
\]
第三行来自弯矩连续，第四行来自剪力跳跃。矩阵 \(\vb H_{XY}(\eta)\) 的行列式恰为两端齐次子空间在源点处的 Wronski 型组合，当 \(\eta\in(\xi_0,1)\) 时非零，故系数唯一可解。
\end{proof}
应用举例：取 $\mathrm{CC}$ 组合（内外缘均固支），内缘基 $\{u_1^{\mathrm C},u_2^{\mathrm C}\}$ 与外缘基 $\{v_1^{\mathrm C},v_2^{\mathrm C}\}$ 均已满足各自固支条件，Green 函数在源点处由四个系数唯一确定，所得挠度即圆环板在环向线载荷下的轴对称响应。
\begin{proof}[解]
轴对称方程 \(\mathcal L_0 f=1\) 中，\(\mathcal L_0=(\dd_\xi^2+\xi^{-1}\dd_\xi)^2\)。对 \(\xi^4/64\) 直接求导，
\[
 \left(\dd_\xi^2+\frac1\xi\dd_\xi\right)\frac{\xi^4}{64}
 =\frac{12\xi^2}{64}+\frac{4\xi^2}{64}
 =\frac{\xi^2}{4},
\]
再作用一次，
\[
 \left(\dd_\xi^2+\frac1\xi\dd_\xi\right)\frac{\xi^2}{4}
 =\frac{2}{4}+\frac{2}{4}
 =1,
\]
故 \(\mathcal L_0(\xi^4/64)=1\)，\(\xi^4/64\) 是一个特解。齐次解空间为 \(\operatorname{span}\{1,\xi^2,\ln\xi,\xi^2\ln\xi\}\)，因此一般解可写成
\[
 f(\xi)=\frac{\xi^4}{64}
 +a_{\mathrm{CF}}\xi^2\ln\xi
 +b_{\mathrm{CF}}\xi^2+c_{\mathrm{CF}}\ln\xi+d_{\mathrm{CF}}.
\]
对这一形式逐项求导并代入 \(\mathcal V[f]=f'''+\xi^{-1}f''-\xi^{-2}f'\)。常数项与 \(\ln\xi\) 项在 \(\mathcal V\) 下分别为 \(0\)；\(\xi^2\) 项给出 \(4\xi\)；\(\xi^2\ln\xi\) 项给出 \(8/\xi\)；特解 \(\xi^4/64\) 给出 \(\xi/2\)。合并得
\[
 \mathcal V[f]
 =\frac{8a_{\mathrm{CF}}}{\xi}+\frac{\xi}{2}
 =\frac{8a_{\mathrm{CF}}+\xi^2}{2\xi}.
\]
内缘自由条件 \(\mathcal V[f](\xi_0)=0\) 给出 \(8a_{\mathrm{CF}}+\xi_0^2=0\)，即
\[
 a_{\mathrm{CF}}=-\frac{\xi_0^2}{8}.
\]
剩余三条件
\[
 \mathcal M[f](\xi_0)=0,
 \qquad f(1)=0,
 \qquad f'(1)=0
\]
是关于 \(b_{\mathrm{CF}},c_{\mathrm{CF}},d_{\mathrm{CF}}\) 的三元线性系统。把 \(a_{\mathrm{CF}}=-\xi_0^2/8\) 代入后显式展开：\(\mathcal M[f]=f''+\nu\xi^{-1}f'\) 在 \(\xi=\xi_0\) 处给出一条方程，\(f(1)=0\) 与 \(f'(1)=0\) 各给出一条方程。具体地，由 \(\ln 1=0\)，\(f(1)=0\) 化为 \(1/64+b_{\mathrm{CF}}+d_{\mathrm{CF}}=0\)，把 \(d_{\mathrm{CF}}\) 用 \(b_{\mathrm{CF}}\) 表出，
\[
 d_{\mathrm{CF}}=-\frac1{64}-b_{\mathrm{CF}},
\]
由 \(f'(1)=0\) 再把 \(c_{\mathrm{CF}}\) 用 \(b_{\mathrm{CF}}\) 表出，最后 \(\mathcal M[f](\xi_0)=0\) 解出 \(b_{\mathrm{CF}}\)。三条方程的系数矩阵行列式为 \(16\Delta\)，\(\Delta=(1-\nu)+(1+\nu)\xi_0^2\)，在 \(\xi_0\in(0,1)\)、\(\nu\in(-1,1)\) 时非零，故解唯一。回代即得正文给出的三个系数。
\end{proof}
\subsection{长方形板的解析解}
\label{sec:phys-sm-rectangular}
继续保持上一节的均匀各向同性材料特例，只把中面区域改为长方形
\[
 \mathcal A=\{(x,y):0<x<a,\ 0<y<b\}.
\]
对长方形板而言，真正能够完全用分离变量处理的情形只有两类。其一是四边简支，此时可以用 Navier 双重正弦级数把本征函数、格林函数和静载响应全部写成显式双重级数。其二是 L\'evy 类边界，即取一对对边，例如 \(x=0,a\)，为简支，而另一对对边 \(y=0,b\) 各自从 \(\mathrm C\)、\(\mathrm S\)、\(\mathrm F\) 三类边界条件中任意选取。这样共有九种有序边界组合；它们的本征值问题都能化成一维常微分方程与一个 \(4\times4\) 特征行列式，因此同样属于解析可解的范围。除此之外的一般四边边界条件虽然仍可数值求解，但不再具有同样简单的分离变量结构，故这里不再展开。
四边简支时，边界条件写成
\[
 w|_{x=0,a}=0,
 \qquad
 w|_{y=0,b}=0,
 \qquad
 M_{nn}|_{x=0,a}=0,
 \qquad
 M_{nn}|_{y=0,b}=0.
\]
由于边界本身已经强制正弦基，最自然的本征函数就是
\[
 \Psi_{m,s}^{\mathrm{SSSS}}(x,y)=\sin\frac{m\pi x}{a}\,\sin\frac{s\pi y}{b},
 \qquad m,s=1,2,\dots.
\]
它们满足
\[
 \nabla_t^2\Psi_{m,s}^{\mathrm{SSSS}}=-\gamma_{m,s}^2\Psi_{m,s}^{\mathrm{SSSS}},
 \qquad
 \gamma_{m,s}^2=\left(\frac{m\pi}{a}\right)^2+\left(\frac{s\pi}{b}\right)^2,
\]
从而
\[
 \nabla_t^4\Psi_{m,s}^{\mathrm{SSSS}}=\gamma_{m,s}^4\Psi_{m,s}^{\mathrm{SSSS}},
 \qquad
 \omega_{m,s}^{\mathrm{SSSS}}=\gamma_{m,s}^2\sqrt{\frac{D}{\rho_{\rm m} h}}.
\]
若模态质量按
\[
 \mathcal M_{m,s}^{\mathrm{SSSS}}=\int_{\mathcal A}\rho_{\rm m} h\,\bigl(\Psi_{m,s}^{\mathrm{SSSS}}\bigr)^2\,\dd a=\frac{\rho_{\rm m} hab}{4}
\]
归一，则四边简支板的静力格林函数便可直接写成显式双重级数
\[
 G^{\mathrm{SSSS}}(x,y;\xi,\eta)=\frac{4}{abD}\sum_{m=1}^{\infty}\sum_{s=1}^{\infty}
 \frac{\sin\dfrac{m\pi x}{a}\sin\dfrac{m\pi \xi}{a}\sin\dfrac{s\pi y}{b}\sin\dfrac{s\pi \eta}{b}}{\left[\left(\dfrac{m\pi}{a}\right)^2+\left(\dfrac{s\pi}{b}\right)^2\right]^2}.
\]
因此任意静载 \(f_\perp(x,y)\) 下的挠度都是
\[
 w(x,y)=\int_0^a\int_0^b G^{\mathrm{SSSS}}(x,y;\xi,\eta)f_\perp(\xi,\eta)\,\dd \xi\,\dd \eta.
\]
若载荷是作用在 \((\xi_*,\eta_*)\) 的点力 \(P\)，即 \(f_\perp=P\delta(x-\xi_*)\delta(y-\eta_*)\)，则得到
\[
 w(x,y)=P\,G^{\mathrm{SSSS}}(x,y;\xi_*,\eta_*).
\]
若载荷是均匀横向负载 \(f_\perp=f_0\)，则只有奇数项留下，直接积分可得经典 Navier 级数
\[
 w(x,y)=\frac{16f_0}{D\pi^6}\sum_{\substack{m,s=1\\ m,s\text{ odd}}}^{\infty}
 \frac{\sin\dfrac{m\pi x}{a}\sin\dfrac{s\pi y}{b}}{ms\left(\dfrac{m^2}{a^2}+\dfrac{s^2}{b^2}\right)^2}.
\]
从本征函数、格林函数到点载与均布载荷的全部公式，在四边简支情形下都已经完全显式了。
固定 \(x=0,a\) 为简支边后，\(x\) 向正弦基自动满足这对边界条件：
\[
 \Psi(x,y)=\sum_{m=1}^{\infty}Y_m(y)\sin(\alpha_m x),
 \qquad \alpha_m=\frac{m\pi}{a}.
\]
于是二维双调和本征值问题对每个 \(m\) 约化为同一个一维四阶算子。对时间因子 \(e^{-\ii\omega t}\)，令
\[
 k^4=\frac{\rho_{\rm m} h\omega^2}{D},
\]
则
\begin{equation}
 \left(\frac{\dd^2}{\dd y^2}-\alpha_m^2\right)^2Y_m=k^4Y_m.
 \label{eq:phys-sm-levy-dynamic-ode}
\end{equation}
这一步体现了 L\'evy 方法的真正内容：几何只通过 \(a,b\) 和 \(\alpha_m\) 进入，另一对边的边界类型只决定同一个四阶算子的自伴实现。
对 \(y=\mathrm{const}\) 的边，去掉不影响齐次条件的公共因子后，可用四个边界迹
\[
 \mathcal W_m[Y]=Y,
 \qquad
 \Theta_m[Y]=Y',
\]
\[
 \mathcal M_m[Y]=Y''-\nu\alpha_m^2Y,
 \qquad
 \mathcal V_m[Y]=Y'''-(2-\nu)\alpha_m^2Y'
\]
统一表示固支、简支和自由边界：
\[
 \mathrm C:(\mathcal W_m,\Theta_m)=(0,0),\qquad
 \mathrm S:(\mathcal W_m,\mathcal M_m)=(0,0),\qquad
 \mathrm F:(\mathcal M_m,\mathcal V_m)=(0,0).
\]
令 \(\boldsymbol\varphi_m(y;k)\) 是\eqnrefs{eq:phys-sm-levy-dynamic-ode}的任意基本解行向量，并把下边 \(X\) 与上边 \(Y\) 的四个边界条件作用在该基本解上，得到 \(4\times4\) 特征矩阵 \(\vb Z_{m,\mathrm L}^{XY}(k)\)。非零本征函数存在当且仅当
\begin{equation}
 \det\vb Z_{m,\mathrm L}^{XY}(k)=0.
 \label{eq:phys-sm-levy-characteristic}
\end{equation}
若其正根记为 \(k_{m,s}^{XY}\)，则
\[
 \omega_{m,s}^{XY}
 =(k_{m,s}^{XY})^2\sqrt{\frac{D}{\rho_{\rm m} h}},
 \qquad
 \Psi_{m,s}^{XY}(x,y)
 =\sin(\alpha_m x)Y_{m,s}^{XY}(y).
\]
因此九种 \(XY\in\{\mathrm C,\mathrm S,\mathrm F\}^2\) 并不是九套理论，只是同一四阶算子的九种边界实现。四边简支对应 \(XY=\mathrm{SS}\)，并进一步退化为前面的 Navier 双正弦谱。
动态本征值问题已经由边界行列式 $\det Z=0$ 解决。静力 Green 核取 $k=0$，用同一一维算子构造显式响应函数。
静力时 \(k=0\)，第 \(m\) 个 Fourier 分量满足
\begin{equation}
 \mathscr L_mY_m
 :=\left(\frac{\dd^2}{\dd y^2}-\alpha_m^2\right)^2Y_m
 =\frac{f_{\perp,m}(y)}{D},
 \qquad
 f_{\perp,m}(y)=\frac{2}{a}\int_0^a f_\perp(x,y)\sin(\alpha_mx)\dd x.
 \label{eq:phys-sm-levy-static-ode}
\end{equation}
相应 Green 核由
\begin{equation}
 \mathscr L_m g_m^{XY}(y,\eta)=\delta(y-\eta)
 \label{eq:phys-sm-levy-green-ode}
\end{equation}
以及两端的 \(XY\) 边界条件唯一决定。由于 \(\mathscr L_m\) 的最高导数系数为一、源只含 \(\delta\) 而不含其导数，源点条件是
\[
[g_m]_{\eta}=0,
\quad[g_m']_{\eta}=0,
\quad[g_m'']_{\eta}=0,
\quad[g_m''']_{\eta}=1.
\]
这与“位移、斜率、弯矩连续而剪力跳跃”的板论表述完全等价。
令
\[
 \boldsymbol\phi_m(y)
 =\bigl(\cosh\alpha_my,\ \sinh\alpha_my,\ y\cosh\alpha_my,\ y\sinh\alpha_my\bigr)
\]
为 \(\ker\mathscr L_m\) 的一组基。则
\[
 g_m^{XY}(y,\eta)=
 \begin{cases}
 \boldsymbol\phi_m(y)\vb c_{m,-}^{XY}(\eta),&y<\eta,\\[2pt]
 \boldsymbol\phi_m(y)\vb c_{m,+}^{XY}(\eta),&y>\eta,
 \end{cases}
\]
其中八个系数由“四个端点条件 + 四个源点匹配条件”组成的 \(8\times8\) 线性系统确定。二维 Green 核因此统一为
\begin{equation}
 G^{XY}(x,y;\xi,\eta)
 =\frac{2}{aD}\sum_{m=1}^{\infty}
 g_m^{XY}(y,\eta)
 \sin(\alpha_mx)\sin(\alpha_m\xi).
 \label{eq:phys-sm-levy-2d-green}
\end{equation}
对任意静载，\(w=G^{XY}*f_\perp\)；点力则直接取对应源点的 Green 核。这里的解析性来自一对简支边所提供的正弦分解，而不是来自某个特定的载荷形状。
均布横向载荷 $f_\perp=f_0$ 是上述 Green 核最直接的统一特例。
\begin{exbox}[L\'evy 类板的均匀横向负载]
令 \(f_\perp(x,y)=f_0\)。由\eqnrefs{eq:phys-sm-levy-static-ode}，
\[
 f_{\perp,m}=\frac{2f_0}{a}\int_0^a\sin\frac{m\pi x}{a}\dd x
 =\begin{cases}
 \dfrac{4f_0}{m\pi},&m\ \text{为奇数},\\[4pt]
 0,&m\ \text{为偶数}.
 \end{cases}
\]
所以只有奇数 \(m\) 被激发。对每个奇数 \(m\)，常数
\[
 Y_m^{(\mathrm p)}=\frac{4f_0}{Dm\pi\alpha_m^4}
\]
是一个特解，因为 \(\mathscr L_m1=\alpha_m^4\)。于是
\begin{equation}
 Y_m^{XY}(y)
 =\frac{4f_0}{Dm\pi\alpha_m^4}
 +\boldsymbol\phi_m(y)\vb d_m^{XY},
 \qquad m\ \text{为奇数},
 \label{eq:phys-sm-levy-uniform-mode}
\end{equation}
而 \(\vb d_m^{XY}\) 只需用两端的四个边界条件确定。最终得到单重级数
\[
 w^{XY}(x,y)
 =\sum_{\substack{m=1\\m\,\mathrm{odd}}}^{\infty}
 Y_m^{XY}(y)\sin\frac{m\pi x}{a}.
\]
当 \(XY=\mathrm{SS}\) 时，这一表示可再化为四边简支的 Navier 双重级数；其他边界组合只改变边界代数系统，而不改变一维母算子 \(\mathscr L_m\)。
\end{exbox}
\begin{proof}[解]
将\eqnrefs{eq:phys-sm-levy-dynamic-ode}因式分解：
\[
\left(\frac{\dd^2}{\dd y^2}-\alpha_m^2-k^2\right)
\left(\frac{\dd^2}{\dd y^2}-\alpha_m^2+k^2\right)Y_m=0.
\]
定义
\[
 p_m^2=\alpha_m^2+k^2,
 \qquad \chi_m^2=k^2-\alpha_m^2,
\]
其中当 \(k<\alpha_m\) 时 \(\chi_m\) 为纯虚数，下面的三角函数按解析延拓理解；这不会改变四维解空间。临界点 \(k=\alpha_m\) 时 \(\chi_m=0\)，应把第四列按非奇异重标度 \(\sin(\chi_my)/\chi_m\to y\) 取极限，而不能把 \(\sin(\chi_my)\) 直接当成零列。其所有边界迹也必须按同一列重标度取极限。除这一可去退化外，一组基本解为
\[
\boldsymbol\varphi_m
=(\cosh p_my,\ \sinh p_my,\ \cos \chi_my,\ \sin \chi_my).
\]
对它逐项作用四个边界迹，得到
\[
\begin{aligned}
\vb R_W&=(\cosh p_my,\ \sinh p_my,\ \cos \chi_my,\ \sin \chi_my),\\
\vb R_{\theta}&=(p_m\sinh p_my,\ p_m\cosh p_my,\ -\chi_m\sin \chi_my,\ \chi_m\cos \chi_my),
\end{aligned}
\]
\[
\begin{aligned}
\vb R_M=(& (p_m^2-\nu\alpha_m^2)\cosh p_my,
(p_m^2-\nu\alpha_m^2)\sinh p_my,\\
&-(\chi_m^2+\nu\alpha_m^2)\cos \chi_my,
-(\chi_m^2+\nu\alpha_m^2)\sin \chi_my),
\end{aligned}
\]
\[
\begin{aligned}
\vb R_V=(&p_m[p_m^2-(2-\nu)\alpha_m^2]\sinh p_my,
 p_m[p_m^2-(2-\nu)\alpha_m^2]\cosh p_my,\\
&\chi_m[\chi_m^2+(2-\nu)\alpha_m^2]\sin \chi_my,
-\chi_m[\chi_m^2+(2-\nu)\alpha_m^2]\cos \chi_my).
\end{aligned}
\]
若 \(\vb c_m=(c_{m,1},c_{m,2},c_{m,3},c_{m,4})^{\mathsf T}\)，则边界量就是相应行向量乘 \(\vb c_m\)。在 \(y=0\) 按 \(X\) 选两行，在 \(y=b\) 按 \(Y\) 选两行，即得正文的 \(\vb Z_{m,\mathrm L}^{XY}(k)\)。四个齐次边界方程存在非零系数向量，当且仅当其行列式为零，从而得到\eqnrefs{eq:phys-sm-levy-characteristic}。
\end{proof}
\begin{proof}[解]
对\eqnrefs{eq:phys-sm-levy-green-ode}，左右两侧分别写成
\[
 g_- =\boldsymbol\phi_m\vb c_-,
 \qquad
 g_+ =\boldsymbol\phi_m\vb c_+.
\]
在 \(y=0\) 对 \(\vb c_-\) 施加下边 \(X\) 的两个条件，在 \(y=b\) 对 \(\vb c_+\) 施加上边 \(Y\) 的两个条件；源点再施加
\[
 g_--g_+=0,\quad
 g_-' -g_+'=0,\quad
 g_-''-g_+''=0,\quad
 g_+'''-g_-'''=1.
\]
前三个条件保证位移、斜率与弯矩在源点连续，第四个则来自 \(\mathscr L_m\) 的最高阶系数为一时 \(\delta\) 源对应的剪力跳跃。把八个方程按 \((\vb c_-,\vb c_+)\) 的分量排成矩阵 \(\vb M_m^{XY}(\eta)\)，其行列式恰为两个二维齐次子空间在源点处的 Wronski 型组合，当 \(\eta\in(0,b)\) 时非零。这八个方程组成唯一的 \(8\times8\) 匹配系统，并给出正文\eqnrefs{eq:phys-sm-levy-2d-green}。具体地，\(\vb M_m^{XY}\) 可按 \(\eta\) 分块为
\[
 \vb M_m^{XY}(\eta)=
 \begin{pmatrix}
  \vb B_X(0) & \vb 0\\
  \vb 0 & \vb B_Y(b)\\
  \boldsymbol\phi_m(\eta) & -\boldsymbol\phi_m(\eta)\\
  \boldsymbol\phi_m'(\eta) & -\boldsymbol\phi_m'(\eta)\\
  \boldsymbol\phi_m''(\eta) & -\boldsymbol\phi_m''(\eta)\\
  -\boldsymbol\phi_m'''(\eta) & \boldsymbol\phi_m'''(\eta)
 \end{pmatrix},
\]
其中 \(\vb B_X,\vb B_Y\) 是两端边界算子对基向量 \(\boldsymbol\phi_m\) 的作用矩阵。
对均布载荷取常数特解 \(Y_m^{(\mathrm p)}=Y_{m,0}\)，其中
\[
 Y_{m,0}=\frac{4f_0}{Dm\pi\alpha_m^4}.
\]
该特解的选取依据是 \(\mathscr L_m1=\alpha_m^4\)，故常数 \(Y_{m,0}\) 恰使 \(\mathscr L_mY_{m,0}=4f_0/(Dm\pi)\) 等于强迫项。于是
\[
\mathcal W_m[Y_{m,0}]=Y_{m,0},
\quad
\Theta_m[Y_{m,0}]=0,
\quad
\mathcal M_m[Y_{m,0}]=-\nu\alpha_m^2Y_{m,0},
\quad
\mathcal V_m[Y_{m,0}]=0.
\]
这里 \(\mathcal W_m,\Theta_m,\mathcal M_m,\mathcal V_m\) 分别为位移、斜率、弯矩与有效剪力算子。因此齐次修正项在一条 \(\mathrm C,\mathrm S,\mathrm F\) 边上分别需要抵消的残量为
\[
\vb r_{\mathrm C}=(-Y_{m,0},0)^{\mathsf T},\qquad
\vb r_{\mathrm S}=(-Y_{m,0},\nu\alpha_m^2Y_{m,0})^{\mathsf T},\qquad
\vb r_{\mathrm F}=(\nu\alpha_m^2Y_{m,0},0)^{\mathsf T}.
\]
以 \(\mathrm C\) 边为例：固支要求位移与斜率同时为零，故残量第一分量为 \(-Y_{m,0}\)（抵消特解位移），第二分量为 \(0\)（特解斜率已为零）。\(\mathrm S\) 边需同时抵消位移与弯矩 \(-\nu\alpha_m^2Y_{m,0}\)；\(\mathrm F\) 边则需抵消弯矩，而特解剪力已为零。把下、上两边相应的两个残量块叠成四维右端，就得到 \(\vb d_m^{XY}\) 的 \(4\times4\) 边界系统
\[
 \vb Z_{m,\mathrm L}^{XY}(k)\,\vb d_m^{XY}
 =\begin{pmatrix}\vb r_X\\\vb r_Y\end{pmatrix},
\]
代回\eqnrefs{eq:phys-sm-levy-uniform-mode}即可。这一处理对任意 \(XY\in\{\mathrm C,\mathrm S,\mathrm F\}^2\) 都适用，只是 \(\vb Z_{m,\mathrm L}^{XY}\) 的行选取不同。
\end{proof}
四边简支与 L\'evy 类边界因此都只是前面一般 Kirchhoff--Love 双调和算子的几何特例：前者把两个方向都谱对角化，后者只对一个方向谱对角化，再用一维自伴边界问题处理另一个方向。具体几何没有改变变分结构、Green 跳跃或谱理论，只改变了实现这些一般结构所用的基函数。
\subsection{Chladni 节线}
\label{sec:phys-sm-chladni-general}
前两节分别把环板和长方形的谱显式写出。本节回到本章标题中的核心物理对象：Chladni 图形。1787 年 Chladni 用琴弓拉振金属板并在其上撒沙，观察到沙粒聚集到振动为零的节线上，形成对称而优美的图案。这些图案并非额外模型，而是双调和算子本征函数的\emph{节点集}。本节目的不是罗列具体图形，而是给出适用于任意有界区域 \(\mathcal A\subset\mathbb R^2\) 和任意自伴边界实现 \(X\in\{\mathrm C,\mathrm S,\mathrm F\}\) 的一般理论，再把环板和长方形作为特例代入。
设 \(\{\Psi_j,\omega_j\}_{j=1}^\infty\) 是双调和算子在区域 \(\mathcal A\) 上对应边界实现 \(X\) 的本征系，按质量内积
\[
 \langle f,g\rangle_\mu=\int_{\mathcal A}\rho_{\rm m} h\,f^*g\,\dd a
\]
归一。第 \(j\) 个 Chladni 图形定义为\emph{节点集}
\begin{equation}
 \mathcal N_j
 :=\{\vb r\in\mathcal A:\Psi_j(\vb r)=0\}.
 \label{eq:phys-sm-chladni-nodal-set}
\end{equation}
在实验中，板面撒布的细沙在驻波腹点处被垂直加速度 \(-\omega_j^2\Psi_j\) 抛起、在节点处静止，因此沙粒沿 \(\mathcal N_j\) 聚集。Chladni 图形因此完全由谱理论决定，不需要在板方程之外引入任何新物理。
节点集的普适结构由以下几条给出。
(i) Courant 节域定理。把节点集 \(\mathcal N_j\) 在 \(\mathcal A\) 中截出的连通开区域称为\emph{节域}，其数目记为 \(\mu(\Psi_j)\)。将本征值按重数排序 \(\omega_1\le\omega_2\le\cdots\)，则
\begin{equation}
 \mu(\Psi_j)\le j.
 \label{eq:phys-sm-courant-nodal-domain}
\end{equation}
这一结论对任意自伴椭圆算子成立，不限于双调和或特定边界。它说明第 \(j\) 个模态最多把板面分成 \(j\) 块振动相位互不连通的区域。对基态 \(j=1\)，\(\mu\le1\) 意味着基态无内节线——板面整体同相振动。
(ii) 节线光滑性与交叉。对双调和算子 \(\nabla_t^4\)，本征函数 \(\Psi_j\) 是实解析的（因为 \(\nabla_t^4\Psi_j=\omega_j^2(\rho_{\rm m}h/D)\Psi_j\) 右端光滑，由椭圆正则性逐次提升可得任意阶可导，再用解析椭圆估计得到实解析）。因此 \(\mathcal N_j\) 是 \(\Psi_j\) 的零水平集。由隐函数定理，在梯度非零处 \(\nabla_t\Psi_j\neq0\)，节线是光滑曲线。在临界点 \(\nabla_t\Psi_j=0\) 处，\(\Psi_j\) 的 Taylor 展开主导项为二次型，节线在该处可形成交叉或孤立点。
(iii) 双调和与 Laplace 节线的区别。对 Laplace 算子 \(-\nabla_t^2\)，节线在内部不相交（Pleijel 定理的推论），两条节线只会在边界相遇。但对双调和算子 \(\nabla_t^4\)，节线\emph{可以在内部交叉}。这是因为 \(\nabla_t^4=(\nabla_t^2)^2\)，若 \(\Psi\) 同时是 Laplace 本征函数 \(-\nabla_t^2\Psi=\lambda\Psi\)，则它自动是双调和本征函数 \(\nabla_t^4\Psi=\lambda^2\Psi\)，而两个 Laplace 本征函数的线性组合 \(\Psi=\Psi_a+c\,\Psi_b\)（若 \(\lambda_a=\lambda_b\) 即简并）的零集正是两条节线的交叉图案。方形板的经典十字形 Chladni 图就是这一机制的产物。
(iv) 边界行为。在边界 \(\partial\mathcal A\) 上，节线与边界的交角受边界条件约束。对固支边 \(\mathrm C\)，\(\Psi_j|_{\partial\mathcal A}=0\) 意味着整条边界都是节点，节线从边界出发的方向由 \(\partial_n\Psi_j=0\) 进一步约束。对简支边 \(\mathrm S\)，\(\Psi_j|_{\partial\mathcal A}=0\) 同样使边界属于节点集，但 \(\partial_n\Psi_j\) 自由，节线以非零角度与边界相遇。对自由边 \(\mathrm F\)，边界一般不属于节点集，节线只在特定点与边界相遇，这些点由 \(M_{nn}=0\) 与 \(V_n=0\) 的额外约束决定。
\begin{proof}[解]
策略：用节域上的 Rayleigh 商构造 \(j-1\) 个试探函数，由极小极大原理推出矛盾。
设 \(\Psi_j\) 的节域为 \(\Omega_1,\ldots,\Omega_\mu\)，其中 \(\mu=\mu(\Psi_j)\)。在每个节域 \(\Omega_\ell\) 上，\(\Psi_j|_{\Omega_\ell}\) 满足 Dirichlet 边界条件（因为在 \(\partial\Omega_\ell\) 上 \(\Psi_j=0\)），且
\[
 \nabla_t^4\Psi_j=\omega_j^2\frac{\rho_{\rm m}h}{D}\Psi_j
 \quad\text{在 }\Omega_\ell\text{ 内}.
\]
因此 \(\omega_j^2\) 是双调和算子在 \(\Omega_\ell\) 上带 Dirichlet 边界的本征值。由 Rayleigh 商，
\[
 \omega_j^2
 =\frac{D\int_{\Omega_\ell}|\nabla_t^2\Psi_j|^2\dd a}
 {\int_{\Omega_\ell}\rho_{\rm m}h\,|\Psi_j|^2\dd a}
 \ge\lambda_1(\Omega_\ell),
\]
其中 \(\lambda_1(\Omega_\ell)\) 是 \(\Omega_\ell\) 上带 Dirichlet 边界的双调和算子第一本征值。
现在构造 \(\mu-1\) 个试探函数。在每个节域上令
\[
 \phi_\ell(\vb r)=
 \begin{cases}
  \displaystyle\frac{\Psi_j|_{\Omega_\ell}}{\|\Psi_j\|_{L^2(\Omega_\ell)}},& \vb r\in\Omega_\ell,\\[6pt]
  0,& \vb r\notin\Omega_\ell.
 \end{cases}
\]
每个 \(\phi_\ell\) 在 \(\mathcal A\) 上 \(H_0^2\)（因为在节域边界上函数值和法向导数都连续为零），且彼此正交：
\[
 \int_{\mathcal A}\rho_{\rm m}h\,\phi_\ell^*\phi_m\dd a=0,
 \qquad \ell\neq m.
\]
取这 \(\mu-1\) 个函数张成的 \((\mu-1)\) 维子空间 \(V_{\mu-1}=\operatorname{span}\{\phi_1,\ldots,\phi_{\mu-1}\}\)。对任意 \(\phi\in V_{\mu-1}\)，其 Rayleigh 商满足
\[
 \mathcal R[\phi]
 =\frac{D\int_{\mathcal A}|\nabla_t^2\phi|^2\dd a}
 {\int_{\mathcal A}\rho_{\rm m}h\,|\phi|^2\dd a}
 =\frac{D\sum_{\ell=1}^{\mu-1}|c_\ell|^2\int_{\Omega_\ell}|\nabla_t^2\Psi_j|^2\dd a/\|\Psi_j\|_{L^2(\Omega_\ell)}^2}
 {\sum_{\ell=1}^{\mu-1}|c_\ell|^2}
 \le\omega_j^2,
\]
因为每个节域上的 Rayleigh 商都等于 \(\omega_j^2\)（\(\Psi_j\) 在每个节域上都是本征函数）。由极小极大原理的第 \((\mu-1)\) 个本征值刻画，
\[
 \omega_{\mu}^2\le\sup_{\phi\in V_{\mu-1}\setminus\{0\}}\mathcal R[\phi]\le\omega_j^2.
\]
由于本征值按递增排列，\(\omega_\mu\le\omega_j\) 蕴含 \(\mu\le j\)。若 \(\mu>j\)，则 \(\omega_\mu>\omega_j\)（严格不等号来自本征值排序），矛盾。故 \(\mu(\Psi_j)\le j\)，即\eqnrefs{eq:phys-sm-courant-nodal-domain}。
\end{proof}
节点集的普适结构确定以后，区域对称群便可用来分类 Chladni 图形。圆板和方形板给出两组可直接计算的实例。
Chladni 图形的丰富程度直接由本征值的简并结构决定。设区域 \(\mathcal A\) 的对称群为 \(G_{\mathcal A}\subset O(2)\)，则双调和算子与 \(G_{\mathcal A}\) 交换，本征空间自动携带 \(G_{\mathcal A}\) 的表示。
(i) 连续对称——圆板。圆板的对称群 \(G_{\mathcal A}=O(2)\) 连续。对 \(m\ge1\)，角向表示 \(e^{\pm\ii m\theta}\) 是二维不可约表示，因此每个径向量子数 \(s\) 对应一个二重简并本征值 \(\omega_{m,s}\)。实本征函数
\[
 \Psi_{m,s}^{(\cos)}=\mathcal R_{m,s}(r)\cos m\theta,
 \qquad
 \Psi_{m,s}^{(\sin)}=\mathcal R_{m,s}(r)\sin m\theta
\]
的节线分别为 \(m\) 条等角节径和 \(m\) 条旋转 \(\pi/(2m)\) 的节径，加上由 \(\mathcal R_{m,s}\) 的径向零点给出的节圆。任意实线性组合
\[
 \Psi_{m,s}^{(\alpha)}=\mathcal R_{m,s}(r)\cos m(\theta-\alpha)
\]
只是把节径整体旋转角度 \(\alpha\)。因此在理想圆对称下，Chladni 图形不是一个离散集合而是一个连续族，参数 \(\alpha\in[0,\pi/m)\)。实验中板的微小不对称（质量分布偏差、驱动点位置）会破缺这一连续简并，选定特定 \(\alpha\)。对 \(m=0\)，无角向简并，节线只有节圆。
(ii) 离散对称——方形板。方形板 \(\mathcal A=(0,a)\times(0,a)\) 的对称群为二面体群 \(D_4\)（四阶旋转和反射，共 8 个元素）。四边简支时本征函数
\[
 \Psi_{m,s}^{\mathrm{SSSS}}=\sin\frac{m\pi x}{a}\sin\frac{s\pi y}{a}
\]
的本征值为
\[
 \omega_{m,s}^2=\frac{D}{\rho_{\rm m}h}\left[\left(\frac{m\pi}{a}\right)^2+\left(\frac{s\pi}{a}\right)^2\right]^2.
\]
当 \(m\neq s\) 时，\(\Psi_{m,s}\) 与 \(\Psi_{s,m}\) 共享同一本征值（\emph{偶然简并}），因为 \(m^2+s^2=s^2+m^2\)。此时任意线性组合
\[
 \Psi=c_1\sin\frac{m\pi x}{a}\sin\frac{s\pi y}{a}
 +c_2\sin\frac{s\pi x}{a}\sin\frac{m\pi y}{a}
\]
都是本征函数，其 Chladni 图随 \((c_1:c_2)\in\mathbb{RP}^1\) 变化。两个极端情形给出截然不同的图案：
\begin{itemize}
 \item \(c_2=0\)：节线为 \(x=k\pi/m\) 或 \(y=l\pi/s\) 的直网格，把板面分成 \(m\times s\) 个矩形节域。
 \item \(c_1=c_2\)：节线方程 \(\sin(m\pi x/a)\sin(s\pi y/a)+\sin(s\pi x/a)\sin(m\pi y/a)=0\)，一般给出弯曲的对角节线，具有 \(D_4\) 对称性。
\end{itemize}
当 \(m=s\) 时无偶然简并，节线为 \(m\times m\) 正交直网格。更高阶的偶然简并（如 \((m,s)=(1,7)\) 与 \((5,5)\) 因 \(1^2+7^2=5^2+5^2=50\) 而简并）会产生更复杂的 Chladni 图，这正是方形板实验中高阶图案丰富多变的数学根源。
(iii) 一般矩形板。矩形板 \(\mathcal A=(0,a)\times(0,b)\)，\(a/b\notin\mathbb Q\) 时，\(m^2/a^2+s^2/b^2\) 几乎不会有非平凡简并（由 Weyl 谱渐近和数论结果，有理比 \(a/b\) 才可能产生偶然简并）。因此一般矩形的 Chladni 图形远不如方形丰富，几乎每个本征值都是单重的，节线为单一模式 \(\sin(m\pi x/a)\sin(s\pi y/b)\) 的直网格。
(iv) 一般区域的 Pleijel 渐近。对任意有界光滑区域，节域数 \(\mu(\Psi_j)\) 的渐近上界由 Pleijel 定理给出：
\begin{equation}
 \limsup_{j\to\infty}\frac{\mu(\Psi_j)}{j}
 \le\frac{4\pi j_{\nu,1}^2}{\operatorname{Area}(\mathcal A)\cdot C_{\mathrm{Weyl}}},
 \label{eq:phys-sm-pleijel-asymptotic}
\end{equation}
其中 \(j_{\nu,1}\) 是相应 Bessel 函数的第一个正零点，\(C_{\mathrm{Weyl}}\) 是 Weyl 谱密度常数。对二维双调和算子，这一常数比 Laplace 情形更小，因此 Pleijel 商严格小于 1，说明 Courant 界 \(\mu\le j\) 在高阶模态上远非紧致——绝大多数高阶模态的节域数显著少于 \(j\)。
上述分类理论在圆板和方形板中有完全显式的实例。
\begin{exbox}[圆板 Chladni 图形的完整分类]
实心圆板 (\(\xi_0\to0\)) 的本征函数为
\[
 \Psi_{m,s}(r,\theta)=\mathcal R_{m,s}(r)\cos m\theta
 \quad\text{或}\quad
 \mathcal R_{m,s}(r)\sin m\theta,
\]
其中径向函数
\[
 \mathcal R_{m,s}(r)=J_m(\beta_{m,s}\xi)I_m(\beta_{m,s}\xi_0)
 -I_m(\beta_{m,s}\xi)J_m(\beta_{m,s}\xi_0)
\]
（\(\xi_0\to0\) 时退化为 \(J_m(\beta_{m,s}\xi)\)，因为 \(I_m(0)=0\) 对 \(m\ge1\)，\(I_0(0)=1\) 对 \(m=0\) 需单独处理）。节线由两个独立因子决定：
角向节径。\(\cos m\theta=0\) 给出 \(m\) 条等角节径 \(\theta=(2k+1)\pi/(2m)\)，\(k=0,\ldots,m-1\)。对 \(\sin m\theta\) 则旋转 \(\pi/(2m)\)。
径向节圆。\(\mathcal R_{m,s}(r)=0\) 的正零点给出 \(s-1\) 个同心节圆（由 Sturm 振荡定理，第 \(s\) 个径向本征函数恰有 \(s-1\) 个内部零点）。
因此第 \((m,s)\) 模态的 Chladni 图形由 \(m\) 条节径和 \(s-1\) 条节圆组成，把圆板分成 \(m\cdot s\) 个节域（满足 Courant 界 \(\mu=m\cdot s\le j\)，其中 \(j\) 是 \((m,s)\) 在排序谱中的位置）。
具体低阶图案：
\begin{itemize}
 \item \((m,s)=(0,1)\)：无节径、无节圆，整板同相振动（基态，\(\mu=1\)）。
 \item \((m,s)=(1,1)\)：1 条节径把圆板分成 2 个半圆 (\(\mu=2\))。
 \item \((m,s)=(2,1)\)：2 条正交节径分成 4 个扇形 (\(\mu=4\))。
 \item \((m,s)=(0,2)\)：1 条节圆分成内圆和外环 (\(\mu=2\))。
 \item \((m,s)=(1,2)\)：1 条节径加 1 条节圆，分成 4 个区域 (\(\mu=4\))。
 \item \((m,s)=(3,1)\)：3 条等角节径分成 6 个扇形 (\(\mu=6\))。
\end{itemize}
\end{exbox}
\begin{exbox}[方形板 SSSS 的 Chladni 图形与偶然简并]
四边简支方形板 (\(a=b\)) 的本征函数 \(\Psi_{m,s}=\sin(m\pi x/a)\sin(s\pi y/a)\) 的节线为直网格
\[
 x=\frac{ka}{m},\quad k=1,\ldots,m-1;
 \qquad
 y=\frac{la}{s},\quad l=1,\ldots,s-1.
\]
节域数为 \(\mu=m\cdot s\)。
无简并情形 （\(m\neq s\) 且 \(m^2+s^2\) 无其他分解）：节线为 \(m\times s\) 矩形直网格。例如 \((1,2)\)：1 条水平节线 \(y=a/2\)，\(\mu=2\)。\((2,3)\)：1 条竖线 \(x=a/2\) 加 2 条横线 \(y=a/3,2a/3\)，\(\mu=6\)。
偶然简并情形 （\(m\neq s\) 且存在 \((p,q)\neq(m,s)\) 使 \(m^2+s^2=p^2+q^2\)）：本征值二重简并，Chladni 图为连续族。最简单的例子是 \((1,7)\) 与 \((5,5)\)：
\[
 1^2+7^2=5^2+5^2=50.
\]
本征函数
\[
 \Psi=c_1\sin\frac{\pi x}{a}\sin\frac{7\pi y}{a}
 +c_2\sin\frac{5\pi x}{a}\sin\frac{5\pi y}{a}
\]
的节线随 \(c_1/c_2\) 从直网格连续变形为对角弯曲节线。\(c_2=0\) 时有 \(1\times7=7\) 个矩形节域；\(c_1=0\) 时有 \(5\times5=25\) 个正方形节域；中间比值给出过渡图案。这正是 Chladni 实验中方形板高阶图案极为丰富的数学原因。
对称组合。取 \(c_1=c_2=1\)，节线方程
\[
 \sin\frac{\pi x}{a}\sin\frac{7\pi y}{a}
 +\sin\frac{5\pi x}{a}\sin\frac{5\pi y}{a}=0
\]
具有 \(D_4\) 全对称性（关于 \(x=y\)、\(x=a-y\)、\(x=a/2\)、\(y=a/2\) 对称），图案呈四重旋转对称。取 \(c_1=1,c_2=-1\) 则关于对角线反对称，节线必过方形中心。
\end{exbox}
Chladni 图形不仅是正问题（给定区域求节线）的输出，还自然引出反问题。
(i) Kac 问题：「能否听出板的形状？」 板的振动谱 \(\{\omega_j\}\) 决定了 Weyl 谱渐近
\[
 N(\omega)=\#\{j:\omega_j\le\omega\}
 \sim\frac{\operatorname{Area}(\mathcal A)}{4\pi}\left(\frac{\rho_{\rm m}h}{D}\right)^{1/2}\omega,
\]
因此从谱可读出面积 \(\operatorname{Area}(\mathcal A)\)。更高阶热迹展开可读出周长、Euler 示性数等。但 Kac 的反问题「谱 \(\{\omega_j\}\to\mathcal A\) 是否一一」的答案是否定的：存在等谱不同构的区域（Milnor 的 16 维反例；二维平面的等谱不 congruent 反例由 Gordon--Wilson--Webb 构造）。因此仅凭 Chladni 图形的频率序列不能唯一反推板的形状。
(ii) 节线反问题。比谱更强的问题是：给定所有 Chladni 图形 \(\{\mathcal N_j\}_{j=1}^\infty\) 的完整集合（不只是频率），能否反推 \(\mathcal A\)？Uhlenbeck 证明了对通有（generic）区域，谱和节点集共同唯一确定算子，从而确定区域。但「通有」排除了等谱反例所在的零测集，因此反问题在可操作意义下的完全回答仍开放。
(iii) Berry 随机波猜想。高阶模态 \(j\gg1\) 的本征函数在局部尺度上趋于高斯随机场（Berry 猜想）。若成立，则高阶 Chladni 图形的节线统计性质（节线长度密度、交叉点密度）由高斯随机场的统计量决定，与具体区域无关。对二维双调和算子，数值实验支持这一猜想在典型区域上成立，但在可积区域（如矩形）上因简并结构而失效——矩形的高阶 Chladni 图形保持直网格，远非随机。
Chladni 原始实验。在金属板下方用琴弓以频率 \(\omega\) 驱动板面振动，板面撒布细沙。沙粒在腹点处被垂直加速度 \(\omega^2|\Psi|\) 抛起、在节线处静止，最终聚集成 \(\mathcal N_j\)。驱动频率连续扫过本征频率 \(\omega_j\) 时，沙图案在共振处突然清晰。这一方法的分辨率受限于沙粒惯性和板阻尼，通常只观察到前几十个模态。
典型实验用黄铜板：密度 \(\rho\approx8.5\times10^3\,\mathrm{kg/m^3}\)，杨氏模量 \(E\approx1.1\times10^{11}\,\mathrm{Pa}\)，泊松比 \(\nu\approx0.34\)，厚度 \(h\approx1.5\,\mathrm{mm}\)。板的抗弯刚度 \(D=Eh^3/[12(1-\nu^2)]\) 决定本征频率的整体尺度：同几何下 \(\omega_j\propto\sqrt{D/(\rho h)}\)，故频率对厚度 \(h\) 极敏感（\(\omega\propto h\)）。
全息干涉法。用激光全息术记录板面驻波的振幅分布 \(|\Psi_j(\vb r)|\)，节线为零等高线。分辨率远高于沙法，可观测数百个模态。
数值方法。给定区域 \(\mathcal A\) 和边界条件 \(X\)，数值生成 Chladni 图形的标准流程为：
\begin{enumerate}
 \item 用有限元或谱方法离散双调和算子，施加边界条件 \(X\)。
 \item 求解广义本征值问题 \(\vb K\vb c=\omega^2\vb M\vb c\)，得到本征对 \((\omega_j,\vb c_j)\)。
 \item 由 \(\vb c_j\) 重建本征函数 \(\Psi_j(\vb r)\)（有限元插值或谱展开）。
 \item 绘制零等高线 \(\Psi_j(\vb r)=0\)，即 Chladni 图形 \(\mathcal N_j\)。
\end{enumerate}
对解析可解区域（圆板、矩形板），步骤 1--3 由前两节的显式公式直接给出，只需绘制零等高线。对一般区域，有限元离散是唯一可行途径，Courant 节域定理可用于验证数值正确性：若数值得到的节域数 \(\mu>j\)，则必有离散或求根错误。
\begin{proof}[解]
Courant 节域定理给出 \(\mu_j\le j\)，但高阶本征函数的节域数远小于 \(j\)。Pleijel 定理给出精确渐近上界。
Pleijel 常数。定义 Pleijel 常数
\begin{equation*}
  \gamma_{\rm Pl}=\frac{4\pi j_0^2}{2\pi}=\frac{2}{j_0^2}\approx0.691,
\end{equation*}
其中 \(j_0\approx2.4048\) 是 \(J_0\) 的第一个正零点。Pleijel 定理（1956）：对有界区域 \(\mathcal D\subset\mathbb R^2\) 上 Dirichlet 拉普拉斯算子，
\begin{equation*}
  \limsup_{j\to\infty}\frac{\mu_j}{j}\le\gamma_{\rm Pl}\approx0.691.
\end{equation*}
因 \(\gamma_{\rm Pl}<1\)，高阶本征函数的节域数严格小于 Courant 上界 \(j\)。
证明骨架。设 \(\mu_j\) 是 \(\phi_j\) 的节域数。每个节域 \(\Omega_k\) 上 \(\phi_j\) 是第一本征函数（Dirichlet，因 \(\phi_j|_{\partial\Omega_k}=0\)），故 \(\lambda_j\ge\lambda_1(\Omega_k)\)。由 Faber--Krahn 不等式（\(\lambda_1(\Omega)\ge\pi j_0^2/|\Omega|\)），\(\lambda_j\ge\pi j_0^2/|\Omega_k|\)，即 \(|\Omega_k|\ge\pi j_0^2/\lambda_j\)。节域互斥：\(\sum_k|\Omega_k|\le|\mathcal D|\)，故
\begin{equation*}
  \mu_j\cdot\frac{\pi j_0^2}{\lambda_j}\le|\mathcal D|,\qquad\mu_j\le\frac{|\mathcal D|\lambda_j}{\pi j_0^2}.
\end{equation*}
由 Weyl 定律 \(j\sim\frac{|\mathcal D|}{4\pi}\lambda_j\)（二维拉普拉斯算子），\(\frac{\mu_j}{j}\le\frac{4\pi}{\pi j_0^2}=\frac{4}{j_0^2}=\gamma_{\rm Pl}\)。
双调和算子的推广。对双调和算子 \(\Delta^2\)（板振动），本征值 \(\lambda_j=\omega_j^2\)，Weyl 定律 \(j\sim\frac{|\mathcal D|}{2\pi}\sqrt{\lambda_j}\)。Faber--Krahn 型不等式对 \(\Delta^2\) 给出 \(\lambda_1(\Omega)\ge C/|\Omega|^2\)（Talenti 型），代入得
\begin{equation*}
  \frac{\mu_j}{j}\le\gamma_{\rm Pl}^{(4)}<1,
\end{equation*}
其中 \(\gamma_{\rm Pl}^{(4)}\) 是四阶 Pleijel 常数（精确值依赖四阶 Faber--Krahn 常数，目前只有数值估计）。因此板的 Chladni 图形在高频时节域数也严格小于 Courant 上界。
\end{proof}
